\chapter*{Conclusion générale et perspectives}
\addcontentsline{toc}{chapter}{Conclusion générale et perspectives}
\markboth{Conclusion générale}{Conclusion générale}

% =============================================================================
\section*{Bilan du travail}
\addcontentsline{toc}{section}{Bilan du travail}
% =============================================================================

Ce projet est parti d'un constat en trois volets. L'interprétation d'un enregistrement
polysomnographique mobilise plusieurs heures d'expertise par patient, ce qui limite
l'accès au diagnostic bien davantage que la capacité d'enregistrement. Cette
interprétation n'est pas parfaitement reproductible, l'accord entre scoreurs
plafonnant autour de \SI{83}{\percent}. Et les modèles d'apprentissage profond qui
pourraient absorber ce volume restent opaques, ce qui interdit au praticien d'engager
sa responsabilité sur une prédiction dont il ignore les motifs.

\hypnoria{} répond à ces trois exigences par une plateforme organisée en deux étages
d'apprentissage. Le premier classe automatiquement les stades du sommeil à partir d'un
enregistrement brut ; le second exploite l'hypnogramme produit pour estimer la sévérité
d'un syndrome d'apnées obstructives, en exposant les facteurs qui déterminent chaque
prédiction. Autour de ces modèles, la plateforme fournit la chaîne applicative attendue
en contexte clinique : dossier patient, archivage sécurisé, annotation, compte rendu
composable, collaboration entre praticiens et traçabilité des actions.

Mené selon le cadre Scrum en trois releases de deux sprints chacune, le projet a livré les
soixante récits de son backlog. Les objectifs fixés au chapitre~\ref{chap:presentation}
ont été atteints dans les proportions que résume le tableau ci-dessous.

\begin{table}[H]
\centering
\caption*{\textbf{Bilan au regard des objectifs}}
\small
\begin{tabular}{@{}L{4.2cm}L{7.6cm}c@{}}
\toprule
\textbf{Objectif} & \textbf{Réalisation} & \textbf{État} \\
\midrule
Automatiser le scorage des stades &
Douze modèles entraînés et quatre ensembles construits ; meilleure exactitude de
\SI{91,66}{\percent} pour un $\kappa$ de \num{0,853} &
atteint \\
\addlinespace[3pt]
Prédire la sévérité du SAOS &
Modèle à quatre classes sur 53 variables ; \SI{57,0}{\percent} d'exactitude et plus de
\SI{95}{\percent} des prédictions à un cran près &
atteint \\
\addlinespace[3pt]
Expliquer chaque prédiction &
Attributions par valeurs de Shapley sur le modèle de sévérité, restituées en facteurs
aggravants et protecteurs &
partiellement \\
\addlinespace[3pt]
Fournir la chaîne applicative &
Trente points d'entrée, deux rôles, archivage cloud, collaboration, journal d'audit,
interface bilingue, pile conteneurisée &
atteint \\
\bottomrule
\end{tabular}
\end{table}

L'objectif d'explicabilité n'est que partiellement atteint : les explications produites
portent sur la prédiction de sévérité, non sur
la classification des stades. Les cartes de saillance et les poids d'attention
initialement envisagés pour les modèles profonds n'ont pas été exposés dans
l'interface. Le mécanisme d'attention existe pourtant dans l'architecture récurrente,
et ses poids sont calculés à chaque inférence : les rendre visibles relève d'un travail
de restitution, non de modélisation.

% =============================================================================
\section*{Apports du projet}
\addcontentsline{toc}{section}{Apports du projet}
% =============================================================================

Sur le plan scientifique, ce travail apporte une comparaison systématique de
trois familles d'architectures sur quatre configurations de montage et de granularité,
là où la littérature compare le plus souvent un modèle à des résultats publiés dans des
conditions différentes. Il apporte surtout deux résultats nuancés que la pratique
courante tend à masquer. D'une part, l'empilement dépasse le meilleur modèle isolé sur
les quatre configurations de classification, mais ne dépasse le vote moyen que sur trois
d'entre elles : la comparaison au vote moyen, rarement rapportée, révèle que l'apport
du méta-apprentissage est réel mais modeste et conditionné au désaccord entre les
modèles combinés. D'autre part, sur la prédiction de sévérité, l'empilement
\emph{dégrade} légèrement la meilleure ligne de base, ce que nous avons rattaché à un
défaut de diversité entre trois modèles tous fondés sur des arbres.

Sur le plan technique, la contribution tient moins aux modèles qu'à
l'intégration. Peu de travaux académiques livrent la chaîne complète permettant à un
praticien d'exploiter un modèle : de la lecture d'un fichier hétérogène jusqu'au compte
rendu archivé, en passant par la gestion des comptes, le cloisonnement des données et
la traçabilité. Deux décisions d'ingénierie ont pesé : l'arbitrage de
résidence des données en mémoire, qui a préservé la qualité de convergence au prix du
volume, et l'archivage en arrière-plan, qui affiche les résultats avant la fin du
transfert.

Sur le plan personnel, ce projet aura été l'occasion de mener un cycle complet,
de la revue de littérature à la conteneurisation, et surtout d'apprendre à
distinguer un résultat d'un artefact de protocole. La découverte, en rédigeant le
chapitre~\ref{chap:release2}, que les matrices de confusion produites par le
pipeline de sévérité portaient des étiquettes permutées, a été plus formatrice que bien
des succès : elle rappelle qu'un chiffre n'a de valeur que si l'on sait exactement ce
qu'il mesure.

% =============================================================================
\section*{Limites}
\addcontentsline{toc}{section}{Limites}
% =============================================================================

Quatre limites bornent la portée de ce travail. Elles ont été analysées en détail dans
les chapitres correspondants ; nous les rassemblons ici.

La première est méthodologique. Les modèles de base ont été évalués sur un
découpage au niveau des époques et non des sujets. Chaque patient contribuant environ un
millier de fenêtres, le modèle a vu, pour presque chaque époque de test, d'autres
époques du même patient et de la même nuit. Les métriques rapportées mesurent donc la
classification d'une époque inconnue chez un patient connu, et surestiment
vraisemblablement la performance en situation clinique. S'y ajoute le fait que les
ensembles ont été évalués selon un protocole différent, ce qui rend les deux séries de
chiffres non comparables entre elles.

La deuxième tient au corpus. Les 196 enregistrements retenus représentent
\SI{3,4}{\percent} de la cohorte disponible, et ils ont été sélectionnés par filtrage
sur la complétude du montage plutôt que par tirage aléatoire. La diversité des
morphologies et des conditions d'enregistrement y est moindre que dans la cohorte
complète.

La troisième concerne l'explicabilité, incomplète comme indiqué plus haut, et dont la
nature doit être précisée : les attributions expliquent le comportement du
modèle, non le phénomène physiologique. Une contribution élevée signale une influence
sur la prédiction, jamais un lien de causalité.

La quatrième est applicative. Les routes d'inférence ne sont pas protégées par
le contrôle d'accès, la conformité au règlement européen reste incomplète sur le
chiffrement au repos, le droit à l'effacement et les durées de conservation, et la
validation repose sur une campagne manuelle plutôt que sur une couverture de tests
automatisés.

% =============================================================================
\section*{Perspectives}
\addcontentsline{toc}{section}{Perspectives}
% =============================================================================

Les travaux à mener se hiérarchisent selon qu'ils conditionnent la
validité des résultats, la mise en service ou l'extension du système.

\subsection*{Consolider la validité des résultats}

La ré-évaluation complète sur un découpage par sujets est la priorité.
Il s'agit d'appliquer un protocole unique (tirage aléatoire des sujets, séparation
stricte entre entraînement, validation et test) aux douze modèles de base comme aux
quatre ensembles. Aucune modification d'architecture n'est nécessaire : seul le
découpage change. Les chiffres obtenus seront probablement inférieurs à ceux rapportés
ici, mais ils seront comparables entre eux et transposables à la pratique clinique.

La validation externe sur la cohorte MESA, dont l'accès a été obtenu mais qui
n'a pas été exploitée, constituerait la seconde étape : elle mesurerait la capacité du
système à généraliser à une population et à un protocole d'acquisition différents.

L'élargissement du corpus suit. Il suppose de renoncer à la
résidence intégrale en mémoire, donc de concevoir un chargement par blocs qui n'affame
pas le processeur graphique, par exemple en préchargeant le bloc suivant pendant le
traitement du bloc courant.

\subsection*{Préparer la mise en service}

L'authentification des routes d'inférence est une correction d'une ligne par
route, à appliquer de part et d'autre. La mise en conformité réglementaire
demande davantage : chiffrement applicatif des données au repos, procédure d'effacement
d'un patient et de ses fichiers en une opération, définition et application de durées de
conservation.

La couverture de tests automatisés doit être construite à trois niveaux :
unitaire sur la chaîne de traitement du signal, d'intégration sur les points d'entrée,
et de bout en bout sur le parcours d'analyse. Les trois assertions de sécurité déjà
présentes dans la chaîne d'intégration continue en constituent l'amorce.

La validation clinique, enfin : aucun praticien n'a évalué le système en
conditions réelles. Une étude de concordance entre le scorage automatique et celui d'un
technicien expérimenté, sur un échantillon d'enregistrements, mesurerait ce que les
métriques ne disent pas : la confiance qu'un professionnel accorde à l'outil.

\subsection*{Étendre le système}

Deux extensions relèvent des promesses initiales du cahier des charges et restent
pertinentes.

Le lissage contextuel des transitions exploiterait le fait qu'un hypnogramme
n'est pas une suite d'époques indépendantes : certaines transitions sont
physiologiquement improbables, et une fenêtre glissante sur les probabilités
prédites corrigerait les aberrations isolées. L'implémentation est simple et le gain
attendu porte précisément sur l'indice de fragmentation, qui alimente le modèle de
sévérité.

L'explicabilité des modèles de classification compléterait celle déjà en place
sur la sévérité. Les poids d'attention du modèle récurrent sont calculés à chaque
inférence sans être exploités ; les exposer indiquerait au praticien quels instants de
l'époque ont déterminé la décision. Des cartes de saillance sur le modèle convolutif
apporteraient une information complémentaire. Les réserves formulées dans la
littérature sur la valeur explicative de l'attention~\cite{jain2019,wiegreffe2019} devront
toutefois être rappelées dans l'interface.

Enfin, le passage à l'échelle suppose de séparer le moteur d'inférence de
l'interface applicative, afin de dimensionner indépendamment les ressources de calcul
et celles de service, et de remplacer les migrations au démarrage par un outil de
migration versionnée.

% =============================================================================
\section*{Mot de la fin}
\addcontentsline{toc}{section}{Mot de la fin}
% =============================================================================

Ce projet ne prétend pas remplacer l'expertise humaine. Un modèle qui atteint le
niveau d'accord observé entre deux experts ne fait pas mieux qu'un expert : il fait la
même chose, plus vite et de manière parfaitement reproductible. Dans un domaine où le
temps d'expertise disponible constitue le facteur limitant, le gain est déjà réel.

La transparence compte cependant autant que la performance brute. Un système qui
expose les facteurs de sa décision permet au praticien de la valider, de la contester,
et de conserver la responsabilité du diagnostic. C'est la condition à laquelle un outil
d'aide à la décision peut être employé au chevet du patient.
